\documentclass[11pt]{article}
\usepackage[margin=1in]{geometry}
\usepackage[T1]{fontenc}
\usepackage{lmodern,amsmath,amssymb,amsthm,microtype}
\usepackage[colorlinks=true,urlcolor=blue,citecolor=blue,linkcolor=blue]{hyperref}
\newtheorem{theorem}{Theorem}
\newtheorem{lemma}[theorem]{Lemma}
\newcommand{\SO}{\mathrm{SO}(2)}
\newcommand{\GL}{\mathrm{GL}^{+}(2,\mathbb R)}
\newcommand{\Aut}{\operatorname{Aut}_{\mathrm{hol}}}
\title{A rotational obstruction to prescribed-end Veech realization}
\author{Alec Kriebel\\\small Independent researcher\\\small\href{https://orcid.org/0009-0001-9320-500X}{ORCID: 0009-0001-9320-500X}}
\date{October 14, 2026}
\begin{document}
\maketitle
\begin{abstract}
The universal prescribed-end realization assertion in Valdez's 2014 Oberwolfach Question~4 has a negative answer in its prescribed matrix-group formulation. An explicit counterexample is the three-point Cantor subspace $X=\{0,2/3,1\}$ and the infinite cyclic group generated by $\frac15\left(\begin{smallmatrix}3&-4\\4&3\end{smallmatrix}\right)$. Every element preserves Euclidean norm, but this group cannot even be contained in the Veech group of an infinite-genus translation surface with exactly three ends, all accumulated by genus. We give the compact-core, conformal and exact arithmetic reductions. The obstruction is an application of established results: Aougab--Patel--Vlamis already imply it in 2020/2021, and Artigiani and collaborators explicitly prove the requisite translation-surface finiteness principle in 2023. No new general finiteness theorem or first-priority claim is made.
\end{abstract}

\paragraph{The question and the answer.}
Question~4 of \cite[p.~887]{OWR}, restated as Question~0.2 in \cite{RMV}, asks whether every closed subspace $X$ of the Cantor set and every countable subgroup $G\leq\GL$ without contracting elements can be realized by a tame infinite-genus translation surface $S$ with no planar ends, $\operatorname{Ends}(S)\cong X$, and full matrix Veech group $G$. Here an end is nonplanar when every end-neighborhood has positive genus. Strict contraction means shortening every nonzero Euclidean vector, equivalently operator norm less than one \cite[p.~674]{PSV}. The witness below preserves every vector norm.

\begin{theorem}\label{main}
Set
\[
 X=\{0,2/3,1\},\qquad A=\frac15\begin{pmatrix}3&-4\\4&3\end{pmatrix},\qquad G=\langle A\rangle.
\]
Then $X$ is a nonempty closed subspace of the middle-thirds Cantor set, and $G\leq\GL$ is countably infinite without contracting elements. No tame infinite-genus translation surface $S$ with no planar ends and $\operatorname{Ends}(S)\cong X$ has a Veech group containing $A$. In particular its Veech group cannot equal $G$.
\end{theorem}

We retain the source's actual underlying surface $S$, not its metric completion. Surfaces are connected, orientable, second countable and without boundary. Translation charts give a Riemann surface structure. If finite cone points belong to $S$, their discrete cone angles are $2\pi k$ and a local parameter $w$ has translation coordinate proportional to $w^k$; thus the complex structure extends across them \cite[p.~676]{PSV}. If a convention omits these points, no points are adjoined here. Infinite-angle completion points are never adjoined as ordinary surface points. The flat metric may be incomplete; no area hypothesis is imposed. The argument therefore applies to the tame surfaces in the question without an extra completeness assumption.

\begin{lemma}\label{proper}
For an infinite-genus Riemann surface $R$ and a compact set $K\subset R$, only finitely many $f\in\Aut(R)$ satisfy $f(K)\cap K\ne\varnothing$.
\end{lemma}
\begin{proof}
Uniformization gives $R=\mathbb H/\Lambda$ with $\Lambda\leq\mathrm{PSL}(2,\mathbb R)$ discrete, free and nonabelian. Indeed infinite genus makes $\pi_1(R)$ nonabelian: the sphere cannot have a nontrivial freely acting holomorphic deck group, and a freely acting holomorphic deck group of the plane consists of translations and is abelian.

Let $N$ be the normalizer of $\Lambda$ in $\mathrm{PSL}(2,\mathbb R)$. Lifting automorphisms identifies $\Aut(R)=N/\Lambda$. Choose noncommuting $a,b\in\Lambda$. If distinct nonidentity $n_j\in N$ approached identity, discreteness of $\Lambda$ would force $n_ja n_j^{-1}=a$ and $n_jb n_j^{-1}=b$ eventually. The centralizer of any nonidentity real orientation-preserving M\"obius transformation is abelian (its elliptic, hyperbolic and parabolic normal forms give this directly). Thus $a,b$ would commute, a contradiction. Hence $N$ is discrete. It is also closed: pairwise quotients of a convergent sequence eventually lie in an identity neighborhood meeting $N$ only at identity, so that sequence is eventually constant.

The action of $\mathrm{PSL}(2,\mathbb R)$ on $\mathbb H$ is proper, since a point stabilizer is compact. Its compact-set transporters therefore meet the closed discrete group $N$ in finitely many elements. To descend this statement, cover $K$ by finitely many open $W_j$ with $\overline{W_j}$ compact and contained in evenly covered neighborhoods $V_j$. Choosing one local section $s_j$ on each $V_j$ yields a compact set $L=\bigcup_j s_j(\overline{W_j})\subset\mathbb H$ whose projection contains $K$; this is a finite covering lift, not a global section. If $[n](K)$ meets $K$, suitable $u,v\in L$ satisfy $nu=\lambda v$ for some $\lambda\in\Lambda$. The representative $\lambda^{-1}n$ transports $L$ to meet itself. There are only finitely many such representatives, hence finitely many cosets.
\end{proof}

\begin{lemma}\label{core}
An infinite-genus surface with exactly $n$ ends, all nonplanar, where $3\leq n<\infty$, has a connected compact subsurface $K$ meeting every homeomorphic image of itself.
\end{lemma}
\begin{proof}
By the classification of orientable surfaces, in the form recalled in \cite[Theorem~1.2]{RMV}, the surface is a compact genus-zero core with $n$ boundary circles, each joined to a one-ended infinite-genus arm. Let $K$ be this core and $U_1,\ldots,U_n$ the open components of its complement. Each $U_i$ carries exactly one \emph{ambient end of $S$}; its intrinsic collar end is irrelevant.

If $h(K)\cap K=\varnothing$, connectedness puts $h(K)$ inside one $U_i$. The connected set $K\cup\bigcup_{j\ne i}\overline{U_j}$ is disjoint from $h(K)$: each closed arm meets $K$ at its cuff. It lies in one component of $S\setminus h(K)$ and carries at least $n-1\geq2$ ambient ends. But those components are $h(U_1),\ldots,h(U_n)$, each carrying exactly one ambient end, since a homeomorphism permutes the ends. This is impossible.
\end{proof}

\paragraph{The rotational obstruction.}
Write $D:\operatorname{Aff}^{+}(S)\to\GL$ for the derivative and $\Gamma(S)=D(\operatorname{Aff}^{+}(S))$ for the full matrix Veech group. The \emph{full preimage}
\[
 I(S)=D^{-1}(\SO)\ \leq\ \Aut(S),\qquad
 D(I(S))=\Gamma(S)\cap\SO
\]
consists locally of maps $z\mapsto\lambda z+b$ with $|\lambda|=1$. They are holomorphic on the regular locus and across included finite cone points by removability in local parameters, as are their inverses. These are the affine isometries of \cite[Section~3.1]{ARSVW}. Lemmas~\ref{proper}--\ref{core} show that $I(S)$, and hence its derivative image, is finite whenever $S$ has finitely many $n\geq3$ nonplanar ends. This uses no splitting of $D$, faithful lift of $\Gamma(S)$, or assumption on the translation kernel.

\begin{proof}[Proof of Theorem~\ref{main}]
The three displayed points belong to the Cantor set, and a finite subset is closed. Direct multiplication gives $A^TA=\mathrm{Id}$ and $\det A=1$, so every integral power is a norm-preserving rotation. To prove infinite order exactly, put $B=5A$ and $u_m=\operatorname{tr}(B^m)$. Cayley--Hamilton gives
\[
 B^2-6B+25\mathrm{Id}=0,\qquad
 u_0=2,\quad u_1=6,\quad u_{m+1}=6u_m-25u_{m-1}\quad(m\geq1).
\]
Induction modulo $5$ yields $u_m\equiv1\pmod5$ for every $m\geq1$. If $A^m=\mathrm{Id}$, then $u_m=2\cdot5^m\equiv0\pmod5$, a contradiction. Thus $G$ is infinite cyclic. If $\Gamma(S)$ contained $A$, its finite intersection with $\SO$ would contain the infinite group $G$. This proves the assertion.
\end{proof}

\paragraph{Prior results and scope.}
The finiteness mechanism is established mathematics. Aougab--Patel--Vlamis \cite[Lemmas~4.1--4.2]{APV} prove the nondisplaceable-core and complete-hyperbolic finiteness statements; the needed lemmas already appear in their July~4, 2020 preprint, pp.~22--23. To apply them here, uniformize the \emph{same} Riemann surface $S$. Its complete Poincar\'e metric is preserved by every holomorphic automorphism. Completeness of this new metric is independent of completeness of the original flat metric. Their lemmas therefore give the entire obstruction above, with no remaining bridge hypothesis. Artigiani--Randecker--Sadanand--Valdez--Weitze-Schmith\"usen \cite[v1, Theorem~1.6; v2, Theorem~1.8]{ARSVW} explicitly prove finite affine isometry groups in the presence of a compact nondisplaceable subsurface. Lemmas~\ref{proper}--\ref{core} are included only to make this application checkable.

The inspected primary sources did not establish an earlier literal negative answer to the named prescribed-matrix question; this bounded search does not certify firstness or continued openness. The 2017 paper's positive conjectural discussion \cite[p.~532]{RMV} is historical context. No new construction, general finiteness theorem or classification is claimed. For $n=1,2$ the core inequality fails; finite rotation groups are not excluded. Nor does the result exclude abstract infinite cyclic Veech groups with a different nonorthogonal generator. The separate Jacob's-ladder lattice and finite-area questions are not addressed.

\paragraph{Verification and disclosure.}
The accompanying package contains the editable source, this PDF, all 27 unchanged original PR919 files, independent mathematical and priority audits, and a standard-library verification entry point. Exact arithmetic controls and immutable-byte checks supplement the written proof; finite enumeration and hashes do not prove universal analytic statements or priority. The original record uses one substantive author approach out of five; this audit adds no author proof-search turn. AI tools were used extensively in research, writing, verification and adversarial review. This preprint is unrefereed; AI-assisted checks are not independent human peer review, and no such human review or formal proof-assistant certification is claimed.

\begin{thebibliography}{9}\small
\bibitem{OWR} F.~Valdez (joint work with C.~Ram\'irez Maluendas), Veech groups of infinite type surfaces, in \emph{Flat Surfaces and Dynamics on Moduli Space}, Oberwolfach Reports \textbf{11} (2014), 885--888; Question~4, p.~887. Report DOI: \href{https://doi.org/10.4171/OWR/2014/15}{10.4171/OWR/2014/15}.
\bibitem{RMV} C.~Ram\'irez Maluendas and F.~Valdez, Veech groups of infinite-genus surfaces, \emph{Algebr. Geom. Topol.} \textbf{17} (2017), 529--560. \href{https://doi.org/10.2140/agt.2017.17.529}{doi:10.2140/agt.2017.17.529}.
\bibitem{PSV} P.~Przytycki, G.~Schmith\"usen and F.~Valdez, Veech Groups of Loch Ness Monsters, \emph{Ann. Inst. Fourier} \textbf{61} (2011), 673--687. \href{https://doi.org/10.5802/aif.2625}{doi:10.5802/aif.2625}.
\bibitem{APV} T.~Aougab, P.~Patel and N.~G.~Vlamis, Isometry groups of infinite-genus hyperbolic surfaces, \emph{Math. Ann.} \textbf{381} (2021), 459--498. \href{https://doi.org/10.1007/s00208-021-02164-z}{doi:10.1007/s00208-021-02164-z}; \href{https://arxiv.org/abs/2007.01982v1}{arXiv:2007.01982v1} (July~4, 2020).
\bibitem{ARSVW} M.~Artigiani, A.~Randecker, C.~Sadanand, F.~Valdez and G.~Weitze-Schmith\"usen, Realizing groups as symmetries of infinite translation surfaces, \href{https://arxiv.org/abs/2311.00158v1}{arXiv:2311.00158v1} (October~31, 2023); \href{https://arxiv.org/abs/2311.00158v2}{v2} (February~23, 2026). Listed as forthcoming in \emph{Algebr. Geom. Topol.} on the \href{https://msp.org/soon/coming.php?jpath=agt}{publisher's forthcoming list}, checked October~14, 2026; no final journal text used.
\end{thebibliography}
\end{document}
